\documentclass[11pt]{article}

% =========================
% NSF GRFP Formatting
% =========================
\usepackage[letterpaper,margin=1in]{geometry}
\usepackage{newtxtext} % Times-style font
\usepackage{setspace}
\usepackage{enumitem}
\usepackage{titlesec}

% No page numbers, headers, or footers
\pagestyle{empty}

% Single spacing
\setstretch{1.0}

% Compact but readable section formatting
\titleformat{\section}[runin]
{\normalfont\bfseries\large}
{}
{0pt}
{}
[:]

\titlespacing*{\section}
{0pt}   % left
{8pt}   % space before
{0.5em} % space after heading

\titleformat{\subsection}[runin]
{\normalfont\bfseries}
{}
{0pt}
{}
[:]

\titlespacing*{\subsection}
{0pt}   % left
{8pt}   % space before
{0.5em} % space after heading

% Slightly cleaner paragraph formatting
\setlength{\parindent}{0pt}
\setlength{\parskip}{5pt}

\begin{document}
	
	% ============================================================
	% PERSONAL, RELEVANT BACKGROUND, AND FUTURE GOALS STATEMENT
	% Maximum length: 3 pages
	%
	% Required NSF headings:
	%   Intellectual Merit
	%   Broader Impacts
	%
	% Before submitting:
	% - Keep margins at 1 inch.
	% - Keep font at 11 pt or larger for Times/Computer Modern.
	% - Do not add page numbers, headers, or footers.
	% - Do not include URLs or DOIs in references.
	% - Remove all drafting comments/placeholders.
	% ============================================================
	
	
\begin{center}
	\textbf{\large Personal, Relevant Background, and Future Goals Statement}
\end{center}
	
	
	% ============================================================
	% OPENING / PERSONAL BACKGROUND
	% Possible topics:
	% - How you became interested in your field
	% - Important educational experiences
	% - What motivates your research interests
	% - Challenges, transitions, or formative experiences
	% - How your interests developed over time
	% ============================================================
	
\section*{Introduction}
%	[Opening paragraph: introduce your background, motivation, and the central
%	thread connecting your experiences and future goals.]
%	
%	[Discuss formative academic, research, professional, or personal experiences.
%	Emphasize growth, initiative, curiosity, problem solving, and sustained
%	engagement with STEM.]
%	
%	[Describe how these experiences shaped the scientific questions or problems
%	you now want to pursue.]

	
\section*{Research and Professional Development}
%	[Describe your most important research/professional experiences.]
%	
%	[Focus on:
%	- the question or problem,
%	- what you specifically did,
%	- methods or skills you developed,
%	- results or outcomes,
%	- what you learned,
%	- how the experience changed your research interests.]
	
	
\section*{Future Goals}
%	[Describe the research questions or areas you hope to pursue in graduate
%	school.]
%	
%	[Explain how graduate education will help you develop the knowledge,
%	technical skills, and independence needed for your long-term STEM career.]
%	
%	[Describe your longer-term career goals and the role you hope to play in your
%	field.]
	
	
	% ============================================================
	% REQUIRED HEADING -- DO NOT RENAME
	%
	% Intellectual Merit should demonstrate your potential to advance
	% knowledge in your field.
	%
	% Possible evidence:
	% - Research experience
	% - Technical preparation
	% - Academic growth
	% - Creativity or independent thinking
	% - Problem solving
	% - Intellectual curiosity
	% - Research questions you want to pursue
	% - Preparation for graduate research
	% ============================================================
	
\section*{Intellectual Merit}
%	[Explain how your background and experiences demonstrate your potential to
%	advance knowledge in your field.]
%	
%	[Connect your previous research, technical skills, coursework, and scientific
%	interests to your readiness for graduate study.]
%	
%	[Explain what intellectual contributions you hope to make through your future
%	research.]
	
	
	% ============================================================
	% REQUIRED HEADING -- DO NOT RENAME
	%
	% Broader Impacts should demonstrate your potential to benefit society.
	%
	% Possible areas:
	% - Teaching or mentoring
	% - Outreach
	% - Broadening participation in STEM
	% - Environmental or societal applications
	% - Public communication
	% - Community engagement
	% - Open/reproducible science
	% - Benefits of your future research
	% ============================================================
	
\section*{Broader Impacts}
%	[Describe ways your previous activities have benefited other people,
%	communities, STEM education, or society.]
%	
%	[Discuss mentoring, teaching, outreach, leadership, service, communication,
%	or other relevant activities.]
%	
%	[Explain how your future graduate research and career could benefit society
%	and how you plan to continue contributing to the broader STEM community.]
%	
	
	% ============================================================
	% REFERENCES
	% Include only if you cite literature in the statement.
	%
	% NSF guidance:
	% - Include journal/publication name.
	% - Do NOT include URLs.
	% - Do NOT include DOIs.
	% ============================================================
	
\subsection*{References}
	
	% Example:
	% 1. Author, A. A., and B. B. Author, ``Article Title,''
	%    \textit{Journal Name}, vol. X, pp. XX--XX, Year.
	
%	[Delete this section if no references are cited, or list references here.]
	
	
\end{document}